The next interesting question is of course if the index of $\mathcal{A}$ is finite, and more particularly, when it is zero. One of the standard methods in answering this question is to exploit the homotopy invariance of the Fredholm index. Consider for example a second-order elliptic operator
\begin{equation}\label{eq:elliptic-dirichlet-index-homotopy-operator}
L(\bm{x},D)u=a_{ij}(\bm{x})\frac{\partial^2u}{\partial x_i\partial x_j}+b_i(\bm{x})\frac{\partial u}{\partial x_i}+c(\bm{x})u
\end{equation}
with Dirichlet boundary condition $B(\bm{x},D)u=u$. We assume the matrix $a_{ij}$ is symmetric and positive definite. We may then consider the one-parameter family of operators
\begin{equation}\label{eq:elliptic-dirichlet-laplace-homotopy}
L_t=(1-t)\Delta+tL,\quad B_t=B.
\end{equation}


If $\Omega$ and the coefficients satisfy the relevant smoothness assumptions, then the assumptions of Theorem~\ref{thm:agmon-douglis-nirenberg-elliptic-estimate} apply for every $t\in[0,1]$; hence the Fredholm index for $(L,B)$ is the same as for Laplace\textquotesingle s equation. In Section~\ref{sec:dirichlet-existence-and-uniqueness}, we proved that the problem $\Delta u=f$ with boundary condition $u=0$ has a unique solution $u\in H^1(\Omega)$ for every $f\in H^{-1}(\Omega)$. Using the inverse trace theorem, we can trivially conclude that there is a unique solution $u\in H^1(\Omega)$ of the problem $\Delta u=f$, $u|_{\partial\Omega}=g$ for every $f\in H^{-1}(\Omega)$, $g\in H^{1/2}(\Omega)$. What we would now like to know is that if $f\in L^2(\Omega)$ and $g\in H^{3/2}(\Omega)$, then $u\in H^2(\Omega)$. This is a statement much along the lines of the first assertion of Theorem~\ref{thm:agmon-douglis-nirenberg-elliptic-estimate}, but is not actually implied by Theorem~\ref{thm:agmon-douglis-nirenberg-elliptic-estimate}. The reason is that for the Dirichlet problem of Laplace\textquotesingle s equation, we would choose $s_1=0$, $t_1=2$ and $r_1=-2$, making $M_1=0$ and $X_{M_1}=H^2(\Omega)$. Hence the theorem asserts higher regularity of $H^2$ solutions if the data are appropriate, but not $H^2$ regularity of $H^1$ solutions. Nevertheless, the regularity of weak solutions can be proved along very similar lines as Theorem~\ref{thm:agmon-douglis-nirenberg-elliptic-estimate}, and Agmon, Douglis and Nirenberg actually state such results for scalar elliptic equations. For second-order equations with Dirichlet conditions, see Sections~\ref{sec:elliptic-interior-regularity} and \ref{sec:elliptic-boundary-regularity}.

A natural question is now to ask for a class of problems to which the approach of Section~\ref{sec:dirichlet-existence-and-uniqueness}, based on the Lax--Milgram lemma, can be extended. This will lead us to Agmon\textquotesingle s condition, to be discussed in Subsection~\ref{subsec:agmon-condition-coercive-problems}. The Lax--Milgram lemma will imply existence of a ``weak\textquotesingle\textquotesingle\ solution, and again the regularity of weak solutions has to be addressed before Theorem~\ref{thm:agmon-douglis-nirenberg-elliptic-estimate} is applicable.

Another interesting question is to characterize the orthogonal complement of the range of $\mathcal{A}$; i.e., what conditions must $\bm{f}$ and $\bm{g}$ satisfy so that the problem $\bm{L}\bm{u}=\bm{f}$ with boundary conditions $\bm{B}\bm{u}=\bm{g}$ is solvable? Usually, one can find a $\bm{u}$ satisfying $\bm{B}\bm{u}=\bm{g}$ by an application of the inverse trace theorem (see next subsection, Subsection~\ref{subsec:adjoint-boundary-value-problem}); hence we are reduced to the case $\bm{g}=0$. This leaves us with the question of characterizing those $\bm{v}$ for which $(\bm{v},\bm{L}\bm{u})=0$ for every $\bm{u}$ satisfying $\bm{B}\bm{u}=0$. By formally integrating by parts, one can obtain an elliptic boundary-value problem for $\bm{v}$, known as the adjoint boundary-value problem. We shall study adjoint boundary-value problems for scalar elliptic equations in the next subsection. Of course, a priori $\bm{v}$ will satisfy the adjoint boundary-value problem only in a ``weak\textquotesingle\textquotesingle\ or ``generalized\textquotesingle\textquotesingle\ sense. Hence the regularity of weak solutions becomes again an important issue. In particular, in order to show that the operator $\mathcal{A}$ is Fredholm, one has to show that the nullspace of the adjoint is finite-dimensional. Of course, one has to show this for weak solutions of the adjoint problem, not just for strong solutions. Indeed, it is possible to prove this. If the coefficients are smooth enough, it turns out that weak solutions of the adjoint problem are actually smooth.
